\documentclass[12pt,a4paper]{report}

\usepackage[utf8]{inputenc}
\usepackage[T1]{fontenc}
\usepackage[french]{babel}
\usepackage{geometry}
\usepackage{setspace}
\usepackage{graphicx}
\usepackage{float}
\usepackage{xcolor}
\usepackage{booktabs}
\usepackage{longtable}
\usepackage{array}
\usepackage{hyperref}
\usepackage{enumitem}
\usepackage{listings}
\usepackage{tikz}
\usetikzlibrary{positioning,arrows.meta,shapes.geometric}

\geometry{margin=2.5cm}
\onehalfspacing
\hypersetup{
  colorlinks=true,
  linkcolor=black,
  urlcolor=blue,
  citecolor=black
}

\definecolor{codebg}{RGB}{248,248,248}
\definecolor{codeframe}{RGB}{210,210,210}
\definecolor{brandolive}{RGB}{90,90,64}

\lstdefinelanguage{JavaScript}{
  keywords={async,await,const,let,var,for,of,if,else,return,new,true,false,null,undefined},
  sensitive=true,
  comment=[l]{//},
  morecomment=[s]{/*}{*/},
  morestring=[b]',
  morestring=[b]",
  morestring=[b]`
}

\lstdefinestyle{codestyle}{
  backgroundcolor=\color{codebg},
  frame=single,
  rulecolor=\color{codeframe},
  basicstyle=\ttfamily\small,
  breaklines=true,
  columns=fullflexible,
  showstringspaces=false,
  keywordstyle=\color{blue},
  commentstyle=\color{gray},
  stringstyle=\color{teal}
}
\lstset{style=codestyle}

\newcommand{\projecttitle}{Saveur - Application de recettes de cuisine et planification de repas}
\newcommand{\studentnames}{A completer : Noms et prenoms des etudiants}
\newcommand{\institution}{A completer : Nom de l'etablissement}
\newcommand{\department}{A completer : Departement / Section}
\newcommand{\teacher}{A completer : Nom de l'enseignant}
\newcommand{\submissiondate}{A completer : Date de remise}

\newcommand{\placeholderfigure}[2]{
  \begin{figure}[H]
    \centering
    \fbox{
      \begin{minipage}[c][7cm][c]{0.88\textwidth}
        \centering
        \textbf{Capture a inserer : #1}\\[0.4cm]
        Placer le fichier image dans \texttt{rapport-latex/figures/}.
      \end{minipage}
    }
    \caption{#2}
  \end{figure}
}

\newcommand{\screenshotfigure}[3]{
  \begin{figure}[H]
    \centering
    \IfFileExists{figures/#1}{
      \includegraphics[width=0.92\textwidth]{figures/#1}
    }{
      \fbox{
        \begin{minipage}[c][7cm][c]{0.88\textwidth}
          \centering
          \textbf{Capture manquante : \texttt{figures/#1}}\\[0.4cm]
          #2
        \end{minipage}
      }
    }
    \caption{#2}
    \label{#3}
  \end{figure}
}

\begin{document}

\begin{titlepage}
  \centering
  \vspace*{1cm}
  {\Large \institution\\}
  {\large \department\\[1.5cm]}

  {\Huge\bfseries \projecttitle\\[0.8cm]}
  {\Large Rapport academique de projet d'examen\\[1.2cm]}

  \begin{tabular}{rl}
    \textbf{Cours :} & Developpement JavaScript\\
    \textbf{Cours :} & Bases de donnees NoSQL\\
    \textbf{Projet :} & Projet 5 - Application de recettes de cuisine\\
    \textbf{Etudiants :} & \studentnames\\
    \textbf{Enseignant :} & \teacher\\
    \textbf{Date de remise :} & \submissiondate\\
  \end{tabular}

  \vfill
  {\large Annee academique 2025-2026\\}
\end{titlepage}

\chapter*{Resume}
\addcontentsline{toc}{chapter}{Resume}

Dans le cadre des cours de Developpement JavaScript et de Bases de donnees NoSQL, nous avons realise une application web full-stack intitulee \textit{Saveur}. L'objectif general est de permettre a un utilisateur authentifie de gerer des recettes de cuisine, de consulter un catalogue public, de filtrer les recettes par categorie, temps et difficulte, d'ajouter des recettes aux favoris, de planifier les repas de la semaine et de generer automatiquement une liste de courses a partir du planning. La solution repose sur React et TypeScript pour l'interface, Node.js et Express pour l'API REST, MongoDB et Mongoose pour la persistance, ainsi que bcrypt et JWT pour la securite. Le projet respecte une architecture separee en modele, vue et controleur : schemas Mongoose, controleurs Express, routes API et pages React. Au terme de l'implementation, l'application est fonctionnelle localement, connectee a MongoDB, et couvre les principales exigences techniques et fonctionnelles du sujet.

\tableofcontents
\listoffigures
\listoftables

\chapter{Introduction}

Les cours de Developpement JavaScript et de Bases de donnees NoSQL ont pour objectif de former les etudiants a la conception d'applications web modernes, interactives et connectees a une base de donnees non relationnelle. Le projet d'examen demande la realisation d'une application web complete, composee d'une interface cliente, d'un serveur Node.js/Express, d'une base MongoDB et d'un systeme d'authentification securise.

La problematique principale consiste a transformer un besoin fonctionnel concret en une solution full-stack coherentement structuree. Pour notre groupe, le sujet retenu est l'application de recettes de cuisine. Ce domaine implique plusieurs entites liees : utilisateurs, recettes, categories, ingredients, etapes de preparation, favoris, planning de repas et liste de courses. La difficulte technique reside dans la modelisation NoSQL de ces donnees et dans la mise en place de relations par references et par embarquement.

Nous avons choisi de developper \textit{Saveur}, une application web permettant de creer et consulter des recettes, de les organiser dans un planning hebdomadaire et de produire automatiquement une liste de courses. L'application vise une utilisation simple et directe : l'utilisateur se connecte, parcourt le catalogue, cree ses propres recettes, ajoute des favoris, planifie ses repas puis obtient les ingredients necessaires.

Ce rapport presente d'abord l'analyse des besoins, puis la conception technique de l'application. Il decrit ensuite l'implementation, l'organisation du code, les tests effectues, les difficultes rencontrees et les ameliorations possibles. Les annexes contiennent les captures d'ecran attendues, le guide d'installation et la repartition des taches.

\chapter{Analyse des besoins}

\section{Objectifs fonctionnels}

Les objectifs fonctionnels definissent ce que l'application doit permettre de faire. Ils proviennent du sujet de projet et des ameliorations que nous avons jugees utiles pour rendre l'application coherente.

\begin{itemize}[leftmargin=*]
  \item Permettre l'inscription et la connexion des utilisateurs.
  \item Securiser les mots de passe par hashage avant leur stockage.
  \item Afficher un catalogue de recettes publiques.
  \item Filtrer les recettes par categorie, difficulte et temps maximal.
  \item Creer une recette personnelle avec titre, description, image, temps, difficulte, portions, ingredients et etapes.
  \item Gerer les ingredients avec quantites et unites.
  \item Ajouter ou retirer une recette des favoris.
  \item Consulter les recettes personnelles de l'utilisateur.
  \item Planifier les repas de la semaine par type de repas : petit dejeuner, dejeuner, diner et snack.
  \item Generer une liste de courses intelligente a partir des recettes planifiees.
  \item Afficher un tableau de bord avec des indicateurs : recettes, favoris, repas planifies et articles de courses.
\end{itemize}

\section{Objectifs non fonctionnels}

\begin{itemize}[leftmargin=*]
  \item \textbf{Securite :} les routes sensibles sont protegees par un token JWT, et les mots de passe sont hashes avec bcrypt.
  \item \textbf{Performance :} les requetes REST renvoient uniquement les donnees utiles et les listes sont limitees lorsque necessaire.
  \item \textbf{Ergonomie :} l'interface est organisee autour de pages claires : tableau de bord, catalogue, planning, courses et compte.
  \item \textbf{Maintenabilite :} le code est organise par responsabilite : modeles, controleurs, routes, middleware, pages et services.
  \item \textbf{Compatibilite :} l'application fonctionne dans les navigateurs modernes grace a Vite et React.
  \item \textbf{Reproductibilite :} le fichier \texttt{README.md} et le fichier \texttt{.env.example} permettent de reproduire l'environnement local.
\end{itemize}

\section{Cibles utilisateurs}

L'application prevoit deux types d'utilisateurs :

\begin{itemize}[leftmargin=*]
  \item \textbf{Utilisateur simple :} peut s'inscrire, se connecter, creer des recettes, consulter le catalogue, ajouter des favoris, planifier ses repas et generer sa liste de courses.
  \item \textbf{Administrateur :} le role existe dans le modele utilisateur. Il peut etre utilise pour des evolutions futures, par exemple la moderation des recettes ou la gestion des categories.
\end{itemize}

\chapter{Conception technique}

\section{Architecture logicielle}

L'application suit une architecture trois tiers. Le client React constitue la couche presentation. L'API Express constitue la couche logique metier et expose les routes REST. MongoDB constitue la couche donnees. Cette separation permet de faire evoluer l'interface sans modifier directement les schemas de base, et de tester l'API separement de la vue.

\begin{figure}[H]
  \centering
  \begin{tikzpicture}[
    node distance=2.1cm,
    box/.style={rectangle, rounded corners, draw=brandolive, very thick, align=center, minimum width=4.2cm, minimum height=1.2cm},
    arrow/.style={-{Latex}, thick}
  ]
    \node[box] (client) {Client React\\Pages + Services};
    \node[box, right=of client] (api) {Serveur Express\\Routes + Controleurs};
    \node[box, right=of api] (db) {MongoDB\\Collections NoSQL};
    \draw[arrow] (client) -- node[above]{HTTP/JSON} (api);
    \draw[arrow] (api) -- node[above]{Mongoose} (db);
    \draw[arrow] (api) to[bend left=18] node[below]{JWT} (client);
  \end{tikzpicture}
  \caption{Architecture trois tiers de Saveur}
  \label{fig:architecture}
\end{figure}

\section{Separation des responsabilites}

Le projet respecte le principe modele, vue, controleur sous une forme adaptee a une application React SPA.

\begin{itemize}[leftmargin=*]
  \item \textbf{Modele :} les schemas Mongoose se trouvent dans \texttt{server/models}. Ils decrivent la structure de chaque collection.
  \item \textbf{Controleur :} les fichiers de \texttt{server/controllers} contiennent la logique de traitement des requetes.
  \item \textbf{Routes :} les fichiers de \texttt{server/routes} associent les endpoints HTTP aux fonctions de controleur.
  \item \textbf{Vue :} les pages React dans \texttt{src/pages} affichent les donnees et gerent les interactions utilisateur.
  \item \textbf{Services frontend :} les fichiers \texttt{src/services} isolent les appels API afin que les vues ne manipulent pas directement \texttt{fetch}.
\end{itemize}

\section{Modelisation NoSQL}

La base utilise sept collections principales. Certaines relations sont gerees par references, et certains sous-documents sont embarques pour optimiser la lecture.

\begin{longtable}{p{3cm}p{5cm}p{5cm}}
  \caption{Collections MongoDB utilisees}\\
  \toprule
  \textbf{Collection} & \textbf{Champs principaux} & \textbf{Relations et justification}\\
  \midrule
  \endfirsthead
  \toprule
  \textbf{Collection} & \textbf{Champs principaux} & \textbf{Relations et justification}\\
  \midrule
  \endhead
  \texttt{users} & \texttt{name}, \texttt{email}, \texttt{passwordHash}, \texttt{role}, timestamps & Referencee par les recettes, favoris, plans et listes. Le mot de passe n'est jamais stocke en clair.\\
  \texttt{categories} & \texttt{name}, \texttt{image}, timestamps & Referencee par \texttt{recipes}. La categorie reste reutilisable pour plusieurs recettes.\\
  \texttt{ingredients} & \texttt{name}, \texttt{defaultUnit}, \texttt{category}, \texttt{caloriesPerUnit} & Referencee dans les ingredients embarques des recettes. Elle facilite la normalisation des noms et unites.\\
  \texttt{recipes} & \texttt{title}, \texttt{description}, \texttt{prepTime}, \texttt{cookTime}, \texttt{difficulty}, \texttt{servings}, \texttt{category}, \texttt{user}, \texttt{ingredients}, \texttt{steps} & Reference vers l'utilisateur et la categorie. Les ingredients et etapes sont embarques car ils sont lus avec la recette.\\
  \texttt{mealplans} & \texttt{user}, \texttt{date}, \texttt{mealType}, \texttt{recipe} & Relie un utilisateur a une recette pour un jour et un type de repas. Index unique sur utilisateur, date et type.\\
  \texttt{favorites} & \texttt{user}, \texttt{recipe}, timestamps & Collection de jonction entre utilisateurs et recettes. Index unique pour eviter les doublons.\\
  \texttt{shoppinglists} & \texttt{user}, \texttt{weekStart}, \texttt{items[]} & Les articles sont embarques car ils representent le resultat calcule d'une liste hebdomadaire.\\
  \bottomrule
\end{longtable}

\begin{figure}[H]
  \centering
  \begin{tikzpicture}[
    node distance=1.5cm and 1.4cm,
    col/.style={rectangle, draw=brandolive, rounded corners, align=center, minimum width=3.1cm, minimum height=0.95cm},
    rel/.style={-{Latex}, thick}
  ]
    \node[col] (users) {users};
    \node[col, right=of users] (recipes) {recipes};
    \node[col, above=of recipes] (categories) {categories};
    \node[col, below=of recipes] (ingredients) {ingredients};
    \node[col, right=of recipes] (favorites) {favorites};
    \node[col, below=of favorites] (mealplans) {mealplans};
    \node[col, below=of mealplans] (shopping) {shoppinglists};

    \draw[rel] (users) -- (recipes);
    \draw[rel] (categories) -- (recipes);
    \draw[rel] (ingredients) -- (recipes);
    \draw[rel] (users) -- (favorites);
    \draw[rel] (recipes) -- (favorites);
    \draw[rel] (users) -- (mealplans);
    \draw[rel] (recipes) -- (mealplans);
    \draw[rel] (mealplans) -- (shopping);
    \draw[rel] (users) to[bend right=20] (shopping);
  \end{tikzpicture}
  \caption{Diagramme simplifie des collections NoSQL}
  \label{fig:collections}
\end{figure}

\section{Authentification}

L'authentification utilise un mecanisme JWT. Lors de l'inscription, le mot de passe est valide puis hashe avec bcrypt. Lors de la connexion, le mot de passe fourni est compare au hash stocke. Si les informations sont correctes, le serveur genere un token JWT contenant l'identifiant de l'utilisateur et son role. Ce token est ensuite envoye par le frontend dans l'en-tete \texttt{Authorization}.

\begin{figure}[H]
  \centering
  \begin{tikzpicture}[
    node distance=1.4cm,
    authstep/.style={rectangle, rounded corners, draw=brandolive, align=center, minimum width=8cm, minimum height=0.8cm},
    arrow/.style={-{Latex}, thick}
  ]
    \node[authstep] (a) {Inscription : nom, email, mot de passe};
    \node[authstep, below=of a] (b) {Hashage bcrypt du mot de passe};
    \node[authstep, below=of b] (c) {Enregistrement dans la collection users};
    \node[authstep, below=of c] (d) {Connexion : verification email + mot de passe};
    \node[authstep, below=of d] (e) {Generation du token JWT};
    \node[authstep, below=of e] (f) {Acces aux routes protegees via middleware};
    \draw[arrow] (a) -- (b);
    \draw[arrow] (b) -- (c);
    \draw[arrow] (c) -- (d);
    \draw[arrow] (d) -- (e);
    \draw[arrow] (e) -- (f);
  \end{tikzpicture}
  \caption{Flux d'authentification}
  \label{fig:authflow}
\end{figure}

\begin{table}[H]
  \centering
  \caption{Routes protegees et autorisations}
  \begin{tabular}{p{5cm}p{4cm}p{5cm}}
    \toprule
    \textbf{Route} & \textbf{Protection} & \textbf{Autorisation}\\
    \midrule
    \texttt{GET /api/auth/me} & JWT requis & Utilisateur connecte\\
    \texttt{POST /api/recipes} & JWT requis & Utilisateur connecte\\
    \texttt{GET /api/recipes/mine} & JWT requis & Proprietaire des recettes\\
    \texttt{DELETE /api/recipes/:id} & JWT requis & Proprietaire ou administrateur\\
    \texttt{GET /api/favorites} & JWT requis & Utilisateur connecte\\
    \texttt{POST /api/favorites/:id/toggle} & JWT requis & Utilisateur connecte\\
    \texttt{GET /api/meal-plans} & JWT requis & Utilisateur connecte\\
    \texttt{POST /api/meal-plans} & JWT requis & Utilisateur connecte\\
    \texttt{GET /api/groceries} & JWT requis & Utilisateur connecte\\
    \texttt{GET /api/dashboard} & JWT requis & Utilisateur connecte\\
    \bottomrule
  \end{tabular}
\end{table}

\chapter{Implementation}

\section{Stack technologique}

\begin{itemize}[leftmargin=*]
  \item \textbf{React 19} pour l'interface utilisateur.
  \item \textbf{TypeScript} pour typer les donnees cote frontend.
  \item \textbf{Vite} pour le serveur de developpement et la generation du build.
  \item \textbf{Tailwind CSS} pour le style visuel.
  \item \textbf{Node.js et Express} pour le backend REST.
  \item \textbf{MongoDB} comme base de donnees NoSQL.
  \item \textbf{Mongoose} pour la definition des schemas et les requetes.
  \item \textbf{bcryptjs} pour le hashage des mots de passe.
  \item \textbf{jsonwebtoken} pour les sessions JWT.
  \item \textbf{lucide-react} pour les icones de l'interface.
\end{itemize}

\section{Organisation du code}

\begin{lstlisting}[language=bash,caption={Arborescence simplifiee du projet}]
server/
  config/          # connexion MongoDB
  controllers/     # logique de traitement des requetes
  data/            # donnees de seed
  middleware/      # authentification et erreurs
  models/          # schemas Mongoose
  routes/          # endpoints Express
  utils/           # helpers JWT
src/
  components/      # composants partages
  pages/           # vues React
  services/        # appels API frontend
  types.ts         # types TypeScript partages
\end{lstlisting}

\section{Fonctionnalites cles}

\subsection{Route API protegee}

La route de creation de recette utilise le middleware \texttt{requireAuth}. Seul un utilisateur connecte peut creer une recette.

\begin{lstlisting}[language=JavaScript,caption={Route protegee pour creer une recette}]
recipeRoutes.post('/', requireAuth, createRecipe);
\end{lstlisting}

\subsection{Hashage du mot de passe}

Lors de l'inscription, le mot de passe est transforme en hash avant insertion dans MongoDB.

\begin{lstlisting}[language=JavaScript,caption={Hashage bcrypt lors de l'inscription}]
const passwordHash = await bcrypt.hash(password, 12);
const user = await User.create({
  name,
  email,
  passwordHash,
  role: 'user',
});
\end{lstlisting}

\subsection{Generation de la liste de courses}

La liste de courses regroupe les ingredients des recettes planifiees pendant la semaine. Les quantites identiques sont additionnees par nom et unite.

\begin{lstlisting}[language=JavaScript,caption={Regroupement des ingredients de la semaine}]
for (const entry of entries) {
  for (const ingredient of entry.recipe?.ingredients || []) {
    const key = `${ingredient.name.toLowerCase()}|${ingredient.unit.toLowerCase()}`;
    const current = itemMap.get(key) || {
      name: ingredient.name,
      quantity: 0,
      unit: ingredient.unit,
      category: 'Meal plan',
      checked: false,
    };

    current.quantity += Number(ingredient.quantity) || 0;
    itemMap.set(key, current);
  }
}
\end{lstlisting}

\subsection{Appel API cote frontend}

Les pages React n'appellent pas directement \texttt{fetch}. Elles utilisent les services dans \texttt{src/services}, ce qui separe la vue et l'acces aux donnees.

\begin{lstlisting}[language=JavaScript,caption={Service frontend pour filtrer les recettes}]
export const getFilteredRecipes = async (filters) => {
  const params = new URLSearchParams();
  if (filters.categoryId) params.set('categoryId', filters.categoryId);
  if (filters.difficulty) params.set('difficulty', String(filters.difficulty));
  if (filters.maxTime) params.set('maxTime', String(filters.maxTime));

  const query = params.toString() ? `?${params}` : '';
  const response = await apiFetch(`/recipes${query}`);
  return response.data;
};
\end{lstlisting}

\chapter{Tests et validation}

\section{Environnement de test}

Les tests ont ete effectues sur un environnement local Windows. MongoDB Server 8.3.4 fonctionne comme service local. L'API Express ecoute sur le port \texttt{5000}. L'interface Vite peut etre lancee sur le port \texttt{3000} ou \texttt{3001}. Les tests ont ete verifies dans un navigateur moderne et par des appels HTTP sur les endpoints principaux.

\section{Plan de tests}

\begin{longtable}{p{1cm}p{5cm}p{5cm}p{2cm}}
  \caption{Cas de test fonctionnels}\\
  \toprule
  \textbf{No} & \textbf{Cas de test} & \textbf{Resultat attendu} & \textbf{Etat}\\
  \midrule
  \endfirsthead
  \toprule
  \textbf{No} & \textbf{Cas de test} & \textbf{Resultat attendu} & \textbf{Etat}\\
  \midrule
  \endhead
  1 & Inscription avec nom, email et mot de passe valide & Compte cree, mot de passe hashe, token retourne & OK\\
  2 & Connexion avec compte demo & Token JWT et donnees utilisateur retournes & OK\\
  3 & Connexion avec mauvais mot de passe & Erreur 401 & OK\\
  4 & Consultation du catalogue public & Liste des recettes publiques affichee & OK\\
  5 & Filtre par categorie & Seules les recettes de la categorie sont affichees & OK\\
  6 & Filtre par difficulte & Recettes limitees au niveau choisi & OK\\
  7 & Creation d'une recette connectee & Recette creee avec ingredients et etapes & OK\\
  8 & Ajout/retrait d'un favori & L'etat favori est modifie en base & OK\\
  9 & Ajout d'une recette au planning & Repas enregistre pour une date et un type & OK\\
  10 & Generation de la liste de courses & Ingredients des recettes planifiees regroupes & OK\\
  11 & Acces a une route protegee sans token & Erreur d'authentification & OK\\
  12 & Build de production & Generation du dossier \texttt{dist} sans erreur & OK\\
  \bottomrule
\end{longtable}

\section{Resultats}

Les tests essentiels valident le fonctionnement de la chaine complete : frontend, API REST, authentification et base MongoDB. Les commandes \texttt{npm run lint} et \texttt{npm run build} ont ete executees avec succes. Le backend retourne egalement un statut positif sur \texttt{/api/health}. Le compte demo \texttt{chef@saveur.local} permet de verifier rapidement le parcours utilisateur.

\section{Difficultes rencontrees et solutions}

\subsection{Migration depuis un prototype mobile/Firebase}

Le projet initial contenait des fichiers Android et Firebase. Or, le sujet exigeait une application web JavaScript avec Node.js, Express et MongoDB. Nous avons donc supprime les fichiers non pertinents et remplace la logique Firestore par une API REST connectee a MongoDB.

\subsection{Modelisation des ingredients}

La generation d'une liste de courses exige de connaitre les quantites et les unites. Une simple liste de noms d'ingredients aurait ete insuffisante. La solution retenue consiste a embarquer dans chaque recette les ingredients avec \texttt{name}, \texttt{quantity}, \texttt{unit} et \texttt{optional}, tout en gardant une collection \texttt{ingredients} reutilisable.

\subsection{Securisation des routes privees}

Les routes telles que les favoris, les recettes personnelles, le planning et les courses doivent appartenir a l'utilisateur connecte. Nous avons ajoute un middleware \texttt{requireAuth} qui verifie le JWT et charge l'utilisateur avant de continuer la requete.

\subsection{Installation locale de MongoDB}

L'installation de MongoDB Server a d'abord echoue a cause d'un manque d'espace disque. Apres liberation d'espace, l'installation a reussi et le service Windows \texttt{MongoDB} a ete demarre. L'application utilise ensuite l'URI locale \texttt{mongodb://127.0.0.1:27017/saveur}.

\chapter{Ameliorations possibles}

\begin{itemize}[leftmargin=*]
  \item Ajouter une interface administrateur complete pour gerer les categories, moderer les recettes et consulter les statistiques globales.
  \item Ajouter l'edition complete des recettes existantes, avec modification des ingredients et des etapes.
  \item Deployer l'application sur une plateforme cloud avec une base MongoDB Atlas.
  \item Ajouter des tests automatises backend et frontend.
  \item Exporter la liste de courses en PDF ou en fichier imprimable.
  \item Ajouter des informations nutritionnelles plus detaillees.
\end{itemize}

\chapter{Conclusion}

Le projet \textit{Saveur} nous a permis de mettre en pratique les competences attendues dans les deux cours. Nous avons construit une application web full-stack avec une interface dynamique, une API REST, une base MongoDB et un systeme d'authentification securise. Les objectifs principaux du sujet sont atteints : gestion des recettes, catalogue avec filtres, favoris, recettes personnelles, tableau de bord, planification de repas et generation de liste de courses.

Sur le plan technique, le projet nous a permis de mieux comprendre la separation des responsabilites, la modelisation NoSQL et la securite cote serveur. Nous avons egalement appris a adapter une application existante a un cahier des charges different, en retirant les parties non conformes et en reconstruisant une architecture plus appropriee.

Comme perspective, l'application peut encore etre enrichie par un espace administrateur, un deploiement en ligne, des tests automatises et des exports PDF. Neanmoins, la version actuelle constitue une base fonctionnelle et conforme aux exigences principales de l'examen.

\appendix

\chapter{Captures d'ecran}

Les figures suivantes doivent etre remplacees par les captures finales avant la remise. Le modele demande au moins six captures : page d'accueil, connexion, formulaire d'ajout, liste avec filtres, dashboard/graphiques et fonctionnalite amelioree.

\screenshotfigure{dashboard.png}{Tableau de bord de l'application Saveur}{fig:dashboard}
\screenshotfigure{login.png}{Page de connexion et d'inscription}{fig:login}
\screenshotfigure{create-recipe.png}{Formulaire de creation d'une recette}{fig:create}
\screenshotfigure{catalog-filters.png}{Catalogue avec filtres par categorie, temps et difficulte}{fig:catalog}
\screenshotfigure{recipe-detail.png}{Detail d'une recette avec ingredients et etapes}{fig:detail}
\screenshotfigure{meal-planner.png}{Planning hebdomadaire des repas}{fig:planner}
\screenshotfigure{grocery-list.png}{Liste de courses generee automatiquement}{fig:grocery}

\chapter{Guide d'installation et d'execution}

\section*{Prerequis}

\begin{itemize}[leftmargin=*]
  \item Node.js installe.
  \item MongoDB Server installe et demarre.
  \item Un terminal dans le dossier du projet.
\end{itemize}

\section*{Etapes}

\begin{enumerate}[leftmargin=*]
  \item Installer les dependances :
  \begin{lstlisting}[language=bash]
npm install
  \end{lstlisting}
  \item Creer le fichier d'environnement :
  \begin{lstlisting}[language=bash]
cp .env.example .env
  \end{lstlisting}
  Sur Windows PowerShell :
  \begin{lstlisting}[language=bash]
Copy-Item .env.example .env
  \end{lstlisting}
  \item Verifier que \texttt{MONGO\_URI} vaut :
  \begin{lstlisting}[language=bash]
mongodb://127.0.0.1:27017/saveur
  \end{lstlisting}
  \item Demarrer l'API :
  \begin{lstlisting}[language=bash]
npm run api
  \end{lstlisting}
  \item Demarrer le frontend dans un autre terminal :
  \begin{lstlisting}[language=bash]
npm run dev
  \end{lstlisting}
  \item Ouvrir le navigateur sur :
  \begin{lstlisting}[language=bash]
http://localhost:3000
  \end{lstlisting}
\end{enumerate}

\section*{Compte demo}

\begin{itemize}[leftmargin=*]
  \item Email : \texttt{chef@saveur.local}
  \item Mot de passe : \texttt{saveur123}
\end{itemize}

\setcounter{chapter}{3}
\chapter{Repartition des taches dans le groupe}

\begin{table}[H]
  \centering
  \caption{Repartition des contributions}
  \begin{tabular}{p{5cm}p{8cm}}
    \toprule
    \textbf{Nom et prenom} & \textbf{Contribution}\\
    \midrule
    A completer & Analyse du sujet, modelisation NoSQL, backend Express\\
    A completer & Interface React, pages et services frontend\\
    A completer & Tests, documentation, captures d'ecran et rapport\\
    A completer & Integration, correction des erreurs et preparation GitHub\\
    \bottomrule
  \end{tabular}
\end{table}

\begin{thebibliography}{9}
\addcontentsline{toc}{chapter}{References}
\bibitem{node} Documentation officielle Node.js, \url{https://nodejs.org/}.
\bibitem{express} Documentation officielle Express, \url{https://expressjs.com/}.
\bibitem{mongodb} Documentation officielle MongoDB, \url{https://www.mongodb.com/docs/}.
\bibitem{mongoose} Documentation officielle Mongoose, \url{https://mongoosejs.com/}.
\bibitem{react} Documentation officielle React, \url{https://react.dev/}.
\bibitem{vite} Documentation officielle Vite, \url{https://vite.dev/}.
\end{thebibliography}

\end{document}
